\section{USAMO 1972}
        \begin{enumerate}
        	\item (\textit{USAMO 1972}) The symbols $(a,b,\ldots,g)$ and $[a,b,\ldots, g]$ denote the greatest common divisor and least common multiple, respectively, of the positive integers $a,b,\ldots, g$. For example, $(3,6,18)=3$ and $[6,15]=30$. Prove that


 $\frac{[a,b,c]^2}{[a,b][b,c][c,a]}=\frac{(a,b,c)^2}{(a,b)(b,c)(c,a)}$.

 



\textbf{Solution 1}

Consider an arbitrary prime $p$. Let $p^\alpha$, $p^\beta$, and $p^\gamma$ be the greatest powers of $p$ that divide $a$, $b$, and $c$. WLOG let $\alpha \leq \beta \leq \gamma$.

 Examining each factor in the equation, we see that the largest power of $p$ that divides the left hand side is $2\gamma -(\beta+\gamma+\gamma) = -\beta$, and the largest power of $p$ that divides the right hand side is 2\alpha -($\alpha + \beta + \alpha) = -\beta$. Since every prime has the same power in both expressions, the expressions are equal. $\blacksquare$

 



\textbf{Solution 2}

Let $p = (a, b, c)$, $pq = (a, b)$, $pr = (b, c)$, and $ps = (c, a)$. Then it follows that $q, r, s$ are pairwise coprime and $a = pqsa'$, $b = pqrb'$, and $c = prsc'$, with $a', b', c'$ pairwise coprime as well. Then, we wish to show
 \[\frac{(pqrsa'b'c')^2}{(pqrsa'b')(pqrsb'c')(pqrsc'a')} = \frac{p^2}{(pq)(pr)(ps)},\]which can be checked fairly easily.

 



\textbf{Solution 3 (very long, but detailed)}

Note: This solution is more of a "bashy" solution, and is easier to think of than the first 2 solutions.

 Dividing both sides of the equation by $(a,b,c)^2$, and then multiplying both sides by $[a,b][b,c][c,a]$ gives us 
 \[\left(\frac{[a,b,c]}{(a,b,c)}\right)^2 = \frac{[a,b][b,c][c,a]}{(a,b)(b,c)(c,a)}.\]Now, we look at the prime factorisations of $a,b,c$. Let 
 \[a=2^{x_1}\cdot 3^{x_2}\cdot 5^{x_3}\cdots,\]
 \[b=2^{y_1}\cdot 3^{y_2}\cdot 5^{y_3}\cdots,\]
 \[c=2^{z_1}\cdot 3^{z_2}\cdot 5^{z_3}\cdots,\]where all of the $x_1, y_1, z_1$ are nonnegative integers (they could be 0).Thus, we have 
 \[[a,b,c]=2^{\max(x_1,y_1,z_1)} \cdot 3^{\max(x_2,y_2,z_2)} \cdot 5^{\max(x_3,y_3,z_3)} \cdots\]and 
 \[(a,b,c)=2^{\min(x_1,y_1,z_1)} \cdot 3^{\min(x_2,y_2,z_2)} \cdot 5^{\min(x_3,y_3,z_3)} \cdots.\]Now, we see that $\left(\frac{[a,b,c]}{(a,b,c)}\right) = 2^{\max(x_1,y_1,z_1) - \min(x_1,y_1,z_1)} \cdot 3^{\max(x_2,y_2,z_2) - \min(x_2,y_2,z_2)} \cdot 5^{\max(x_3,y_3,z_3) - \min(x_3,y_3,z_3)} \cdots.$ Squaring the LHS will just double all the exponents on the RHS.

 
Also, we have $[a,b]=2^{\max(x_1,y_1)} \cdot 3^{\max(x_2,y_2)} \cdot 5^{\max(x_3,y_3)} \cdots$. We can also build similar equations for $[b,c],[c,a],(a,b),(b,c)(c,a)$, using 2 of $x,y,z$ and either the minimum of the exponents or the maximum.We see that when we multiply $[a,b]$, $[b,c]$, and $[a,c]$ together, the exponents of the prime factorization will be $\max(x_i,y_i)+ \max(y_i,z_i)+ \max(z_i,x_1)$, where $i$ is chosen for the $i$'th prime. When we multiply $(a,b)$, $(b,c)$, and $(a,c)$ together, the exponents of the prime factorization will be $\min(x_i,y_i)+ \min(y_i,z_i)+ \min(z_i,x_1)$, where $i$ is chosen for the $i$'th prime.

 Thus, when we divide the former by the latter, we have $\max(x_i,y_i)+ \max(y_i,z_i)+ \max(x_i,z_i) - (\min(x_i,y_i)+ \min(y_i,z_i)+ \min(x_i,z_i))$. We wish to prove that this is equal to the exponents we got from $([a,b,c]/(a,b,c))^2$.

 In other words, this problem is now down to proving that 
 \[2(\max(x_i,y_i,z_i) - \min(x_i,y_i,z_i)) = \max(x_i,y_i)+ \max(y_i,z_i)+ \max(x_i,z_i) - (\min(x_i,y_i)+ \min(y_i,z_i)+ \min(x_i,z_i))\] for any nonnegative integers $x_i, y_i, z_i$. WLOG, let $x_i \le y_i \le z_i$. Thus, computing maximums and minimums, this equation turns into 
 \[2(z_i - x_i) = y_i + z_i + z_i - x_i - y_i - x_i.\]

 Finally, canceling out the $y_i$'s and collecting like terms, we see the RHS is $2z_i - 2x_i$, which is equal to the LHS. Thus, this equation is true, and we are done. $\square$

 



\textbf{Solution 4}

Let
 \[a=dzxA\]
 \[b=dxyB\]
 \[c=dyzC\]where $\gcd(A,B,C)=\gcd(x,y,z)=1$. (Existence proof missing.) Then, $\operatorname{lcm}(a,b,c)^2=(dxyzABC)^2$ and $\operatorname{lcm}(a,b)=dxyzAB$. Using this cyclically, we have
 \[\text{LHS}=\frac{d^2x^2y^2z^2A^2B^2C^2}{(dxyzAB)(dxyzBC)(dxyzCA)}=\frac{1}{dxyz}.\]Now, note that $\gcd(a,b,c)^2=d^2$ and $\gcd(a,b)=dx$. Using this cyclically, we have
 \[\text{RHS}=\frac{d^2}{(dx)(dy)(dz)}=\frac{1}{dxyz}.\]-brainiacmaniac31

 \textbf{Note} If one thinks of the problem from a set theory vantage point, it is quite intuitive to recognize from PIE:
 \[\text{lcm}(a,b,c) = \frac{a \cdot b \cdot c \cdot \text{gcd}(a,b,c)}{\text{gcd}(a,b) \cdot \text{gcd}(b,c) \cdot \text{gcd}(c,a)}\]where $\text{lcm}(a,b,c)$ can be thought of as the union, $a \cdot b \cdot c$ can be thought of as the raw addition of sets, and the $\text{gcd}$ function can be thought of as intersections. 

 


	\item (\textit{USAMO 1972}) A given tetrahedron $ABCD$ is isosceles, that is, $AB=CD, AC=BD, AD=BC$. Show that the faces of the tetrahedron are acute-angled triangles.


 



\textbf{Solution 1}

Suppose $\triangle ABD$ is fixed.By the equality conditions, it follows that the maximal possible value of $BC$ occurs when the four vertices are coplanar, with $C$ on the opposite side of $\overline{AD}$ as $B$.In this case, the tetrahedron is not actually a tetrahedron, so this maximum isn't actually attainable.

 For the sake of contradiction, suppose $\angle ABD$ is non-acute.Then, $(AD)^2\geq (AB)^2+(BD)^2$.In our optimal case noted above, $ACDB$ is a parallelogram, so
 \begin{align*} 2(BD)^2 + 2(AB)^2 &= (AD)^2 + (CB)^2 \\ &= 2(AD)^2 \\ &\geq 2(BD)^2+2(AB)^2.  \end{align*}
However, as stated, equality cannot be attained, so we get our desired contradiction.

 



\textbf{Solution 2}

It's not hard to see that the four faces are congruent from SSS Congruence. Without loss of generality, assume that $AB\leq BC \leq CA$. Now assume, for the sake of contradiction, that each face is non-acute; that is, right or obtuse. Consider triangles $\triangle ABC$ and $\triangle ABD$. They share side $AB$. Let $k$ and $l$ be the planes passing through $A$ and $B$, respectively, that are perpendicular to side $AB$. We have that triangles $ABC$ and $ABD$ are non-acute, so $C$ and $D$ are not strictly between planes $k$ and $l$. Therefore the length of $CD$ is at least the distance between the planes, which is $AB$. However, if $CD=AB$, then the four points $A$, $B$, $C$, and $D$ are coplanar, and the volume of $ABCD$ would be zero. Therefore $CD>AB$. However, we were given that $CD=AB$ in the problem, which leads to a contradiction. Therefore the faces of the tetrahedron must all be acute.

 



\textbf{Solution 3}

Let $\vec{a} = \overrightarrow{DA}$, $\vec{b} = \overrightarrow{DB}$, and $\vec{c} = \overrightarrow{DC}$. The conditions given translate to
 \begin{align*} \vec{a}\cdot\vec{a} &= \vec{b}\cdot\vec{b} + \vec{c}\cdot\vec{c} - 2(\vec{b}\cdot\vec{c}) \\ \vec{b}\cdot\vec{b} &= \vec{c}\cdot\vec{c} + \vec{a}\cdot\vec{a} - 2(\vec{c}\cdot\vec{a}) \\ \vec{c}\cdot\vec{c} &= \vec{a}\cdot\vec{a} + \vec{b}\cdot\vec{b} - 2(\vec{a}\cdot\vec{b}) \end{align*}
We wish to show that $\vec{a}\cdot\vec{b}$, $\vec{b}\cdot\vec{c}$, and $\vec{c}\cdot\vec{a}$ are all positive. WLOG, $\vec{a}\cdot\vec{a}\geq \vec{b}\cdot\vec{b}, \vec{c}\cdot\vec{c} > 0$, so it immediately follows that $\vec{a}\cdot\vec{b}$ and $\vec{a}\cdot\vec{c}$ are positive. Adding all three equations,
 \[\vec{a}\cdot\vec{a} + \vec{b}\cdot\vec{b} + \vec{c}\cdot\vec{c} = 2(\vec{a}\cdot\vec{b} + \vec{a}\cdot\vec{c} + \vec{b}\cdot\vec{c})\]In addition,
 \begin{align*} (\vec{a} - \vec{b} - \vec{c})\cdot(\vec{a} - \vec{b} - \vec{c})&\geq 0 \\ \vec{a}\cdot\vec{a} + \vec{b}\cdot\vec{b} + \vec{c}\cdot\vec{c}&\geq 2(\vec{a}\cdot\vec{b} + \vec{a}\cdot\vec{c} - \vec{b}\cdot\vec{c}) \\ 2(\vec{a}\cdot\vec{b} + \vec{a}\cdot\vec{c} + \vec{b}\cdot\vec{c})&\geq 2(\vec{a}\cdot\vec{b} + \vec{a}\cdot\vec{c} - \vec{b}\cdot\vec{c}) \\ \vec{b}\cdot\vec{c}&\geq 0 \end{align*}
Equality could only occur if $\vec{a} = \vec{b} + \vec{c}$, which requires the vectors to be coplanar and the original tetrahedron to be degenerate.

 



\textbf{Solution 4}

Suppose for the sake of contradiction that $\angle BAC$ is not acute. Since all three sides of triangles $BAC$ and $CDB$ are congruent, those two triangles are congruent, meaning $\angle BDC=\angle BAC>90^{\circ}$. Construct a sphere with diameter $BC$. Since angles $BAC$ and $BDC$ are both not acute, $A$ and $D$ both lie on or inside the sphere. We seek to make $AD=BC$ to satisfy the conditions of the problem. This can only occur when $AD$ is a diameter of the sphere, since both points lie on or inside the sphere. However, for $AD$ to be a diameter, all four points must be coplanar, as all diameters intersect at the center of the sphere. This would make tetrahedron $ABCD$ degenerate, creating a contradiction. Thus, all angles on a face of an isosceles tetrahedron are acute.

 



\textbf{Solution 5}

Proof by contradiction: Assume at least one of the tetrahedron's faces are obtuse. WLOG, assume $\angle BAC$ is an obtuse angle. Using SSS congruence to prove that all four faces of the tetrahedron are congruent also shows that the angles surrounding point $A$ are congruent to angles in triangle $ABC$; namely, $\angle BAD$ is congruent to $\angle ABC$, and $\angle CAD$ is congruent to $\angle ACB$. Since the internal angles of triangle $ABC$ must add to 180 degrees, so do the angles surrounding point A. Now lay triangle $ABC$ on a flat surface. A diagram would make it clear that from the perspective of an aerial view, the "apparent" measures of $\angle BAD$ and $\angle CAD$ (which are most likely distorted visually, assuming these angles stick up vertically from the flat surface on which triangle $ABC$ lies) can never exceed the true measures of those angles (equality happens when these angles also lie flat on top of triangle $ABC$). This means that these two angles can never join to form side AD (because $\angle BAC$ is more than the sum of $\angle BAD$ and $\angle CAD$ - a direct consequence of the facts that $\angle BAC$ is obtuse and all three angles add up to 180 degrees), so the tetrahedron with obtuse triangle faces is impossible.

 



\textbf{Solution 6}

Lemma: given triangle $ABC$ and the midpoint of $BC$, which we will call $M$, we can say that if $AM > \frac{BC}{2}$, then $\angle A < 90$.Proof: Since $M$ is the midpoint of $BC$, $BM = MC = \frac{BC}{2}$. Since it is given that $AM > \frac{BC}{2}$, we can substitute $\frac{BC}{2}$ to get two inequalities:
 \[AM > CM, \quad AM > BM.\]The above inequalities imply that $\angle C > \angle CAM$ and $\angle B > \angle BAM$. Adding these inequalities and simplifying the RHS, we have that $\angle C + \angle B > \angle A$. Adding $\angle A$ to both sides, replacing the LHS with $180$ and dividing by $2$ gets us that $\angle A < 90$. This is our desired inequality, so we are done.Note that all faces of this tetrahedron are congruent, by SSS.In particular, we will use that $\triangle ABD \cong \triangle BAC$. WLOG, assume that $\angle ADB$ is the largest angle in triangle $ABD$.Because $\triangle ABD \cong \triangle BAC$, the median from $D$ to $AB$ is equal length to the median from $C$ to $AB$. These points meet at $E$, the midpoint of $AB$. By the triangle inequality,
 \[DE + CE > CD.\]By substituting $CD$ with $AB$ (this is a given in the problem) and $CE$ with $DE$, and then dividing by $2$, we get that
 \[DE > \frac{AB}{2}.\]By the lemma we showed at the start, this implies that $\angle ADB < 90$, and since we said that $\angle ADB$ was the largest angle, triangle $ADB$ must be acute. Since all of the faces of this tetrahedron are congruent, then, all of the faces must be acute.

 <i>Alternate solutions are always welcome. If you have a different, elegant solution to this problem, please \href{https://artofproblemsolving.com/wiki/index.php?title=1972\_USAMO\_Problems/Problem\_2&action=edit}{add it} to this page.</i>

 


	\item (\textit{USAMO 1972}) A random number selector can only select one of the nine integers 1, 2, ..., 9, and it makes these selections with equal probability. Determine the probability that after $n$ selections ($n>1$), the product of the $n$ numbers selected will be divisible by 10.


 



\textbf{Solution 1}

For the product to be divisible by 10, there must be a factor of 2 and a factor of 5 in there.

 The probability that there is no 5 is $\left( \frac{8}{9}\right)^n$.

 The probability that there is no 2 is $\left( \frac{5}{9}\right)^n$.

 The probability that there is neither a 2 nor 5 is $\left( \frac{4}{9}\right)^n$, which is included in both previous cases. 

 The only possibility left is getting a 2 and a 5, making the product divisible by 10.By complementarity and principle of inclusion-exclusion, the probability of that is $1- \left( \left( \frac{8}{9}\right)^n + \left( \frac{5}{9}\right)^n - \left( \frac{4}{9}\right)^n\right)=\boxed{1-(8/9)^n-(5/9)^n+(4/9)^n}$.

 



\textbf{Solution 2 (Recursion)}

Define $a_n$ as the probability that the product is divisible by $10$ after selection $n$. Likewise, define $b_n$ and $c_n$ with divisibility by $2$ and $5$, respectively. Define $d_n$ to be the chance of dividing neither $2$ nor $5$ (and thus not $10$ either) similarly.

 It is clear that $d_n=\left(\frac{4}{9}\right)^n$. Now we can define other recursive formulas:

 We are able to reach a $b$ state via selecting a non-$5$ from a $b$ state and selecting an even number from a $d$ state. Thus its formula is $b_n=\frac{8}{9}b_{n-1}+\frac{4}{9}d_{n-1}$.

 We are able to reach a $c$ state via selecting a non-even number from a $c$ state and selecting a $5$ from a $d$ state. Thus its formula is $c_n=\frac{5}{9}c_{n-1}+\frac{1}{9}d_{n-1}$.

 Finally, to reach an $a$ state, we can select a $5$ from a $b$ state and select an even number from a $c$ state. We can also reach $a_n$ from $a_{n-1}$ because of the fact that once the product is divisible by $10$, it will always be divisible by $10$ regardless of the following selections. Thus its formula is $a_n=a_{n-1}+\frac{1}{9}b_{n-1}+\frac{4}{9}c_{n-1}$.

 For our formula for $b_n$, we can substitute to find that $b_n=\frac{8}{9}b_{n-1}+\left(\frac{4}{9}\right)^n$. Solving this recursion yields $b_n=\left(\frac{8}{9}\right)^n-\left(\frac{4}{9}\right)^n$.

 For our formula for $c_n$, we can substitute to find that $c_n=\frac{5}{9}c_{n-1}+\frac{1}{9}\cdot\left(\frac{4}{9}\right)^{n-1}$. Solving this recursion yields $c_n=\left(\frac{5}{9}\right)^n-\left(\frac{4}{9}\right)^n$.

 Finally, substituting the values into the $a_n$ formula yields $a_n=a_{n-1}+\frac{1}{9}\left(\left(\frac{8}{9}\right)^{n-1}+4\cdot\left(\frac{5}{9}\right)^{n-1}-5\cdot\left(\frac{4}{9}\right)^{n-1}\right)$. Solving this recursion yields the final answer, $\boxed{1-\left(\frac{8}{9}\right)^n-\left(\frac{5}{9}\right)^n+\left(\frac{4}{9}\right)^n}$.

 <i>Alternate solutions are always welcome. If you have a different, elegant solution to this problem, please \href{https://artofproblemsolving.com/wiki/index.php?title=1972\_USAMO\_Problems/Problem\_3&action=edit}{add it} to this page.</i>

 


	\item (\textit{USAMO 1972}) Let $R$ denote a non-negative rational number. Determine a fixed set of integers $a,b,c,d,e,f$, such that for \textit{every} choice of $R$, 


 $\left|\frac{aR^2+bR+c}{dR^2+eR+f}-\sqrt[3]{2}\right|<|R-\sqrt[3]{2}|$

 



\textbf{Solution}

Note that when $R$ approaches $\sqrt[3]{2}$, $\frac{aR^2+bR+c}{dR^2+eR+f}$ must also approach $\sqrt[3]{2}$ for the given inequality to hold. Therefore

 
 \[\lim_{R\rightarrow \sqrt[3]{2}} \frac{aR^2+bR+c}{dR^2+eR+f}=\sqrt[3]{2}\]

 which happens if and only if

 
 \[\frac{a\sqrt[3]{4}+b\sqrt[3]{2}+c}{d\sqrt[3]{4}+e\sqrt[3]{2}+f}=\sqrt[3]{2}\]

 We cross multiply to get $a\sqrt[3]{4}+b\sqrt[3]{2}+c=2d+e\sqrt[3]{4}+f\sqrt[3]{2}$. It's not hard to show that, since $a$, $b$, $c$, $d$, $e$, and $f$ are integers, then $a=e$, $b=f$, and $c=2d$.

 Note, however, that this is a necessary but insufficient condition. For example, we must also have $a^2<2bc$ to ensure the function does not have any vertical asymptotes (which would violate the desired property). A simple search shows that $a=0$, $b=2$, and $c=2$ works.

 


	\item (\textit{USAMO 1972}) A given convex pentagon $ABCDE$ has the property that the area of each of the five triangles $ABC$, $BCD$, $CDE$, $DEA$, and $EAB$ is unity. Show that all pentagons with the above property have the same area, and calculate that area. Show, furthermore, that there are infinitely many non-congruent pentagons having the above area property.


 
                    \begin{figure}[H]
                        \centering
                        \includegraphics[width=0.5\textwidth]{Problem/USAMO_I_1972_5_Problem_Figure_1.png}
                    \end{figure}
                    


 



\textbf{Solution}

\textbf{Lemma:} Convex pentagon $A_0A_1A_2A_3A_4$ has the property that $[A_0A_1A_2] = [A_1A_2A_3] = [A_2A_3A_4] = [A_3A_4A_0] = [A_4A_0A_1]$ if and only if $\overline{A_{n - 1}A_{n + 1}}\parallel\overline{A_{n - 2}A_{n + 2}}$ for $n = 0, 1, 2, 3, 4$ (indices taken mod 5).

 \textbf{Proof:} For the "only if" direction, since $[A_0A_1A_2] = [A_1A_2A_3]$, $A_0$ and $A_3$ are equidistant from $\overline{A_1A_2}$, and since the pentagon is convex, $\overline{A_0A_3}\parallel\overline{A_1A_2}$. The other four pairs of parallel lines are established by a symmetrical argument, and the proof for the other direction is just this, but reversed.

 
                    \begin{figure}[H]
                        \centering
                        \includegraphics[width=0.5\textwidth]{Solution/USAMO_I_1972_5_Solution_Figure_1.png}
                    \end{figure}
                    

 Let $A'B'C'D'E'$ be the inner pentagon, labeled so that $A$ and $A'$ are opposite each other, and let $a, b, c, d, e$ be the side lengths of the inner pentagon labeled opposite their corresponding vertices. From the lemma, pentagon $A'B'C'D'E'$ is similar to pentagon $ABCDE$ with $AB = m(A'B')$ and parallelogram $ABCB'$ (and cyclic) has area 2. Supposing
 \begin{align*} P &= [ABCDE] \\ Q &= [A'B'C'D'E'] \\ R &= \sum_{\text{cyc}}[AC'D'] \\ S &= \sum_{\text{cyc}}[ABD'], \end{align*}
we have
 \begin{align*} \sum_{\text{cyc}}[ABCB'] &= 5Q + 3R + 2S \\ 10 &= 3Q + R + 2P. \end{align*}
Since $BCDC'$ and $CDED'$ are parallelograms, $BD' = EC' = (m - 1)a$. Triangle $EB'C'$ is similar to triangle $ECB$, so $me = BC = \left(\frac{2m - 1}{m - 1}\right)e$, and with the requirement that $m > 1$,
 \[m = \frac{2m - 1}{m - 1}\Rightarrow m = \frac{3 + \sqrt{5}}{2}.\]Now, we compute that $[AC'D'] = \frac{1}{2m - 1} = \sqrt{5} - 2$, and similar computation for the other four triangles gives $R = 5\sqrt{5} - 10$. From the aforementioned pentagon similarity, $Q = P/m^2 = \left(\frac{7 - 3\sqrt{5}}{2}\right)P$. Solving for $P$, we have
 \begin{align*} 10 &= 3\left(\frac{7 - 3\sqrt{5}}{2}\right)P + (5\sqrt{5} - 10) + 2P \\ 20 - 5\sqrt{5} &= \left(\frac{25 - 9\sqrt{5}}{2}\right)P \\ P &= \frac{40 - 10\sqrt{5}}{25 - 9\sqrt{5}} \\ &= \frac{5 + \sqrt{5}}{2}. \end{align*}


 To show that there are infinitely many pentagons with the given property, we start with triangle $A'B'C'$ and construct the original pentagon using the required length ratios. This triangle has three continuous degrees of freedom, and the condition that the equal areas be specifically unity means that there are two continuous degrees of freedom in which the configuration can vary, which implies the desired result.

 Alternatively, one can shear a regular pentagon, a transformation which preserves areas and the property of a pair of lines being parallel.

 


\end{enumerate}